\section{人身犯罪 - 侵犯人身自由的犯罪}
\begin{enumerate}
  \item 最重要的。
  \item \zx{非法拘禁}。剥夺他人人身自由。\zd{实力控制被害人}。
        \begin{enumerate}
          \item 拘禁行为跟空间有关系。单位、校园为标准、小区。可以成为 N 多罪的实行行为。
          \item 特殊类型。\al{吊塔上楼后取走梯子、洗澡把人衣服偷拿走}。
          \item 欺骗型拘禁。\al{骗人到电梯，切断电源。待在这里别动，给你找胡歌来}。
          \item \zx{非法拘禁实行行为之内的。致人重伤、死亡} \yjt \bs{结果加重犯}。和介入因素结合考。
                \begin{enumerate}
                  \item 实行行为之内应理解为任何想\zx{继续实力控制被害人人身自由、意志、通信的行为}。
                  \item \al{甲绑架乙，2 日后发现\xhx{被绳子勒死}了}。
                  \item \al{甲为了拘禁乙，大力\xhx{扭其胳膊}}。考过。
                  \item \al{甲非法拘禁乙，大力反扭乙的胳膊，致其胳膊折断}。
                \end{enumerate}
          \item \zx{非法拘禁实行行为之外的。使用暴力（过失）致人伤残/死亡} \yjt \bs{拟制}，\zx{定故意伤害/杀人}。
                \begin{enumerate}
                  \item \al{甲绑架乙，乙辱骂甲，甲鞭子抽打\xhx{不慎}到太阳穴，乙死亡}。
                  \item \al{甲绑架乙，乙辱骂甲，甲\xhx{气愤}之至，杀死乙}。
                \end{enumerate}
        \end{enumerate}
  \item \bs{绑架}（最重要）。以向\bs{第三人}\zx{提出不法要求为目的}（即不限于要钱），\zx{实力控制人质，向第三人提出不法要求}。
        \begin{enumerate}
          \item 保护法益。人质的生命和自由安全，不保护第三人的财产。
          \item 实力控制，不法要求。不局限要钱。典型。非法拘禁 A，敲诈勒索 B。
          \item \zx{行为公式}。
                \begin{enumerate}
                  \item 目的 1 = \zx{拘禁目的}、目的 2 = \zx{勒索财物的目的}（核心） \yjt 没有 \yjt 非法拘禁。
                  \item 行为 1 = \zx{拘禁行为}、Opt(行为 2) = \zx{向 Relate(A, B) 勒索财物的行为}。
                \end{enumerate}
          \item 完整的行为公式。\bs{拘禁目的 + 勒索财物目的 \yjt 拘禁行为 \yjt Opt(向 Relate(A, B) 勒索财物的行为)}。
                \al{甲绑架乙女 \yjt \yi 卖到非洲。\er 向乙丈夫索要钱财。\san 劫个色。\si 把钱拿出来}。
          \item 成立。拘禁目的 + 勒索财物目的 \yjt 实行拘禁行为。
                \begin{enumerate}
                  \item 人质要合格。
                        \begin{enumerate}
                          \item 相对于第三人而言，绑架的人得是特定人。\al{甲绑架乙\xhx{向特朗普}要钱}。
                          \item 行为时足以让一般第三人担忧。\al{甲绑架乙妻，乙说爱杀就杀，我正跟她离婚}。
                        \end{enumerate}
                  \item 有绑到合格人质的危险性、可能性即可。
                        \begin{enumerate}
                          \item \al{甲听闻丙很有钱，校园门口绑架其子，但实际\xhx{丙膝下无子} \yjt 无可能性，\xhx{未遂}}。
                          \item \al{校园门口，甲绑架乙儿向乙要钱，结果\xhx{绑的是丙的儿子} \yjt 有可能性}。
                        \end{enumerate}
                \end{enumerate}
          \item 既遂。成立 + 拘禁行为\zx{把合格人质控制到手}。
                \al{甲绑到乙，让乙通知其母亲\xhx{赎人}，\scx{乙怕母亲心脏病发作，谎称开车撞了人}。后面电话怎么打无所谓的}。
          \item 特殊问题。
                \begin{enumerate}
                  \item \zx{为索取债务非法扣押他人 \yjt 缺少勒索财物的目的 \yjt 定非法拘禁}。含非法债务，\al{赌债}。
                        \begin{enumerate}
                          \item \al{乙欠甲 1W，拒不还钱，甲绑其妻向乙讨债}。
                          \item \al{乙欠甲 1W，拒不还钱，甲绑其表哥向乙讨债}。
                          \item \bs{视情况还要回答出不构成抢劫罪、不构成绑架罪}。
                        \end{enumerate}
                  \item 结合犯。\zd{绑架持续期间 + 杀人 \hz 故意伤害致人重伤、死亡 = 绑架加重处罚}。不包括过失。
                        \begin{enumerate}
                          \item \al{甲绑架乙，没要到钱 \yjt 撕票、打成重伤、回来发现窒息而死}。
                          \item \al{甲绑架乙，为制服乙，对其拳打脚踢}。
                          \item 观点展示。\yi 要求既遂，未遂 \yjt 数罪并罚。\er 只要求杀人行为。
                        \end{enumerate}
                  \item \bs{共犯模型}。\al{甲戊绑架丙}。
                        \begin{enumerate}
                          \item \al{喊来乙\xhx{帮忙做饭}、喂饭 \yjt 乙中立的帮助行为}。
                          \item \al{喊来乙\xhx{帮忙看押} \yjt 乙承继的共犯}。
                          \item \al{\xhx{骗}乙说丙父欠他钱，让乙\xhx{帮忙看押}，自己去讨债 \yjt 甲乙非法拘禁的共同犯罪，甲绑架，\xhx{乙对绑架目的并不知情，不构成绑架}}。
                          \item \al{向其妻勒索被拒绝，甲怕杀不死乙，喊来丙帮忙 \yjt 甲绑架 + 杀人，乙仅故意杀人}。
                          \item \al{向丙妻要 20W，丙妻觉得这是除掉丈夫的好机会，没有给钱也不报警，甲恼羞成怒，杀害乙}。
                          \item \al{向丙妻要 20W，丙妻说我给你 100W，把他杀了}。
                        \end{enumerate}
                \end{enumerate}
          \item \bs{不要和抢劫混}。
        \end{enumerate}
  \item 拐卖类。
        \begin{enumerate}
          \item 拐卖妇女、儿童。实力控制 + 出卖目的。
                \begin{enumerate}
                  \item 自然人。
                  \item 目的 1 = 实力控制妇女、儿童的目的、目的 2 = 出卖目的。
                  \item 行为 1 = 实力控制妇女、儿童的行为、Opt(行为 2) = 出卖行为。
                  \item 成立。实力控制妇女、儿童的目的 + 出卖目的 \yjt 实行实力控制妇女、儿童的行为。
                  \item 实行行为。绑架 \dy 贩卖 \dy 收买 \dy 拐骗。\al{接送、中转}。
                        \begin{enumerate}
                          \item 绑架。\yi 实力控制行为（行为 1） 实施到位，此时不要求有出卖行为（行为 2）。
                          \item 贩卖。没有实力控制行为（行为 1），卖掉既遂。\al{捡到弃婴卖掉}。
                          \item 收买。出卖目的，买到既遂。\al{二手贩子 \dui、买来自己养 \cuo}。
                          \item 拐骗。出卖目的，主要体现在骗，卖掉既遂。\al{盲山盲井}。
                          \item 卖和送的区别。卖 = 有偿对价/行情价。
                                \begin{itemize}
                                  \item \al{代孕问题 \yjt 不算卖自己的孩子}。
                                  \item \al{甲将乙女骗至丙甲，甲对丙说让乙做你老婆，然后收取丙 10 块钱作为车费，坐公交车离去 \yjt 没有出卖目的}。
                                \end{itemize}
                        \end{enumerate}
                  \item 结果加重犯。致人重伤、死亡。
                  \item 结合犯。拐卖 + 强奸 = 拐卖加重处罚。
                \end{enumerate}
          \item 收买被拐卖的妇女、儿童。没有出卖目的。
                \begin{enumerate}
                  \item 既遂。买到手。
                  \item 仙人跳。\al{甲乙合谋，甲说把乙卖给丙，丙付钱后，乙逃走 \yjt 甲乙诈骗。丙对象不能犯，是被害人}。
                \end{enumerate}
          \item 拐骗儿童。没有出卖目的。\al{自己养 \cuo、甲\scx{为报复乙}（动机），将乙的孩子 2000 \xhx{卖}给他人 \dui}。
          \item \sjt 两个罪的罪数。原则犯收买/拐骗罪后犯任何罪都要数罪并罚。
                \begin{enumerate}
                  \item 实务却不是这样，判轻了。\al{徐州铁链女}。
                  \item 例外。收买/拐骗 + 拐卖 = 拐卖。
                \end{enumerate}
        \end{enumerate}
  \item \yi 几个罪名都有拘禁行为，区别是主观目的。\er 注意承继共犯问题。既遂之后参与进来仍然可以构成共犯。
\end{enumerate}
